\documentclass[11pt]{article}

\usepackage[a4paper,margin=2.4cm]{geometry}
\usepackage{amsmath,amssymb}
\usepackage{siunitx}
\usepackage{xcolor}
\usepackage[hidelinks]{hyperref}
\usepackage{enumitem}

\newcommand{\Ri}{R_{\mathrm{i}}}
\newcommand{\Ro}{R_{\mathrm{o}}}
\newcommand{\dvec}{\hat{\mathbf{d}}}

\title{Navigation geometry\\
\large Ray--solid distances and the detector block grid in \texttt{PTCryspMC.jl}}
\author{}
\date{\today}

\begin{document}
\maketitle

\section{Setup and conventions}
\label{sec:setup}

Transport advances a photon by the smaller of the distance to the next interaction
and the distance to the next geometric boundary. This note gives the closed-form
boundary distances for the two solids the simulation uses --- a \emph{cylinder}
(the phantom and the air world) and a \emph{cylindrical shell} (the detector ring)
--- and the arithmetic that maps a point in the shell to a readout block.

Each solid is defined in its own \emph{local frame}, centred at the origin with its
axis along $z$; the physical-volume layer subtracts the placement (and, later, a
rotation) before calling these routines, so here the solid is always origin-centred.
A ray is
\begin{equation}
  \mathbf{r}(t) = \mathbf{p} + t\,\dvec, \qquad t \ge 0,\quad \|\dvec\| = 1 ,
\end{equation}
with $\mathbf{p}=(p_x,p_y,p_z)$ the current point and $\dvec=(d_x,d_y,d_z)$ the
direction. Write the transverse (radial) part as
\begin{equation}
  \rho^2(t) = (p_x + t d_x)^2 + (p_y + t d_y)^2 = a\,t^2 + b\,t + (p_x^2 + p_y^2),
  \qquad a = d_x^2 + d_y^2,\quad b = 2(p_x d_x + p_y d_y).
\end{equation}
A crossing counts only if it lies ahead of the start surface, $t > \varepsilon$
(\texttt{SURFACE\_EPS}, \SI{1}{pm}); a direction component is treated as parallel to
a surface when its magnitude is below \texttt{PARALLEL\_EPS}. Boundary tests are
inclusive ($\le$), so a ray through an exact edge is always claimed by some surface.

\section{Ray--cylinder (convex)}
\label{sec:cyl}

A cylinder of radius $R$ and half-length $H$ occupies $\rho^2 \le R^2$, $|z|\le H$.
A ray meets its boundary on the lateral wall or on an end cap.

\paragraph{Lateral wall.} $\rho^2(t)=R^2$ is the quadratic
\begin{equation}
  a\,t^2 + b\,t + c = 0, \qquad c = p_x^2 + p_y^2 - R^2 ,
\end{equation}
with (for $a>\texttt{PARALLEL\_EPS}$ and discriminant $\Delta=b^2-4ac\ge 0$) roots
\begin{equation}
  t_\pm = \frac{-b \pm \sqrt{\Delta}}{2a}.
\end{equation}
A root is a valid wall crossing if it falls within the axial extent,
$|p_z + t\,d_z| \le H$.

\paragraph{End caps.} The planes $z=\pm H$ are hit at
\begin{equation}
  t = \frac{\pm H - p_z}{d_z} \quad (d_z \neq 0),
\end{equation}
valid if the crossing lies within the disc, $(p_x+t d_x)^2+(p_y+t d_y)^2 \le R^2$.

\paragraph{Entry and exit.} Gather the valid forward crossings and keep the nearest
and farthest, $t_{\mathrm{near}} = \min$, $t_{\mathrm{far}} = \max$. A convex solid
has at most two:
\begin{itemize}[nosep]
\item \textbf{Exit:} $\texttt{distance\_to\_exit} = t_{\mathrm{far}}$ (the surface the
  ray leaves through), or $\infty$ if it never meets the cylinder. For a point
  \emph{inside} the cylinder there is a single forward crossing,
  $t_{\mathrm{near}}=t_{\mathrm{far}}$, which is that exit.
\item \textbf{Entry:} $\texttt{distance\_to\_entry} = t_{\mathrm{near}}$ when
  $t_{\mathrm{near}} < t_{\mathrm{far}}$ (two distinct crossings $\Rightarrow$ the
  start is outside and the ray enters at the near surface), else $\infty$ (already
  inside, or a miss).
\end{itemize}

\section{Ray--cylindrical shell (non-convex)}
\label{sec:shell}

The detector ring is a shell of inner radius $\Ri$, wall thickness $w$, outer radius
$\Ro=\Ri+w$ and half-length $H$:
\begin{equation}
  \Ri^2 \le \rho^2 \le \Ro^2, \qquad |z| \le H .
\end{equation}
The bore ($\rho<\Ri$) makes the shell \emph{non-convex}: a chord can cross the wall,
traverse the empty bore, and cross the wall again. Near/far crossings no longer
suffice; we work with the parameter intervals along the ray instead.

The shell is the outer solid cylinder, clipped to the axial slab, with the inner
bore removed. Each piece is an interval in $t$:
\begin{align}
  \text{outer disc } \rho^2 \le \Ro^2: &\quad
    \mathcal{O} = [\,t_{o-},\,t_{o+}\,]
    \ \text{(roots of } a t^2 + b t + p_\perp^2 - \Ro^2 = 0),\\
  \text{axial slab } |z|\le H: &\quad
    \mathcal{Z} = [\,t_{z-},\,t_{z+}\,], \quad
    t_{z\mp} = \frac{\mp H - p_z}{d_z},\\
  \text{bore } \rho^2 < \Ri^2: &\quad
    \mathcal{B} = (\,t_{i-},\,t_{i+}\,)
    \ \text{(roots of } a t^2 + b t + p_\perp^2 - \Ri^2 = 0),
\end{align}
where $p_\perp^2 = p_x^2+p_y^2$ and the parallel cases are: if $a\le\texttt{PARALLEL\_EPS}$
the radial conditions are constant, so $\mathcal{O}$ is all of $\mathbb{R}$ when
$p_\perp \le \Ro$ (else empty) and $\mathcal{B}$ is empty unless $p_\perp<\Ri$; if
$d_z\approx 0$ then $\mathcal{Z}=\mathbb{R}$ when $|p_z|\le H$ (else empty). The shell
occupancy is
\begin{equation}
  \mathcal{S} = \bigl(\mathcal{O}\cap\mathcal{Z}\bigr)\setminus\mathcal{B}.
\end{equation}
Writing $\mathcal{O}\cap\mathcal{Z}=[t_a,t_b]$ with
$t_a=\max(t_{o-},t_{z-})$, $t_b=\min(t_{o+},t_{z+})$ (empty if $t_a>t_b$), and
removing the open bore $(t_{i-},t_{i+})$ leaves at most two closed intervals:
\begin{equation}
  \mathcal{S} =
  \begin{cases}
    \{[t_a,t_b]\} & \text{bore disjoint from } [t_a,t_b],\\[2pt]
    \{[t_a,t_{i-}],\,[t_{i+},t_b]\} & \text{bore strictly inside } [t_a,t_b],\\[2pt]
    \{[t_{i+},t_b]\}\ \text{or}\ \{[t_a,t_{i-}]\} & \text{bore overlaps one end.}
  \end{cases}
\end{equation}

\paragraph{Entry and exit.} With the photon at $t=0$:
\begin{itemize}[nosep]
\item \textbf{Exit:} if $0$ lies in a shell interval $[s,e]\in\mathcal{S}$ (the photon
  is in the wall), $\texttt{distance\_to\_exit}=e$ --- the end of the current piece,
  whether that end is the outer wall, the inner wall (into the bore) or a cap.
\item \textbf{Entry:} $\texttt{distance\_to\_entry}$ is the smallest interval start
  $s>\varepsilon$ --- the near face of the next wall segment ahead.
\end{itemize}
The two-interval case is exactly the diametral chord: the photon enters the wall at
$t_a$, leaves into the bore at $t_{i-}$, re-enters at $t_{i+}$ and finally exits at
$t_b$.

\section{The $(\varphi,z)$ block / wheel partition}
\label{sec:grid}

The ring is read out as $N_\varphi$ azimuthal \emph{blocks} $\times$ $N_z$ axial
\emph{wheels}. We model this as a structured grid on the shell rather than as
$N_\varphi N_z$ separate volumes: a block is one $(\varphi,z)$ cell spanning the full
wall depth, so the cost of locating a point is independent of the block count and the
dead gaps between crystals are skipped (the wall is treated as continuous).

For a point $(x,y,z)$ in the shell's local frame, the cell indices are pure
arithmetic. With $\Delta\varphi = 2\pi/N_\varphi$ and $\Delta z = 2H/N_z$,
\begin{align}
  \varphi &= \operatorname{atan2}(y,x) \bmod 2\pi \in [0,2\pi), &
  i_\varphi &= \left\lfloor \varphi/\Delta\varphi \right\rfloor \in \{0,\dots,N_\varphi-1\},\\
  \zeta &= z + H \in [0,2H], &
  i_z &= \left\lfloor \zeta/\Delta z \right\rfloor \in \{0,\dots,N_z-1\},
\end{align}
($\varphi=0$ on the $+x$ axis, increasing counter-clockwise; the top edge $z=H$ is
clamped to $i_z=N_z-1$). The linear block id is
\begin{equation}
  \mathrm{block} = i_z\,N_\varphi + i_\varphi \in \{0,\dots,N_\varphi N_z - 1\}.
\end{equation}
The radial depth $\rho-\Ri\in[0,w]$ is the depth of interaction; it characterises the
hit within a block but does not enter the block id.

\paragraph{Match to the physical crystals.} The grid is fixed so that one cell equals
one physical crystal. For CRYSP1M ($\Ri=\SI{38.7}{cm}$, $H=\SI{51.2}{cm}$,
$N_\varphi=48$, $N_z=20$):
\begin{itemize}[nosep]
\item azimuthal pitch at the inner face $\;s_\varphi = \Ri\,\Delta\varphi
  = 2\pi\,(\SI{387}{mm})/48 \approx \SI{50.6}{mm}$;
\item axial pitch $\;\Delta z = 2H/N_z = \SI{1024}{mm}/20 = \SI{51.2}{mm}$.
\end{itemize}
Both reproduce the wrapped \SI{50}{mm} crystal, and $N_\varphi N_z = 48\times20 = 960$
matches the crystal count of the scanner. Overspill --- a photon that Compton-scatters
across a cell boundary --- is then detected by comparing the block ids of its
deposits, with no stepping at the boundaries; the gaps can be reinstated later as thin
non-interacting layers, at which point the $\varphi$- and $z$-plane crossing distances
(the nearest $\varphi=k\,\Delta\varphi$ and $z=-H+k\,\Delta z$ planes ahead) are added
to the boundary set.

\section{Multi-volume navigation}
\label{sec:nav}

\S\ref{sec:cyl}--\ref{sec:shell} give the boundary distance for a photon inside
\emph{one} solid. A photon history crosses several: it is born in the water phantom,
may scatter there, leaves into the air gap, and reaches the crystal ring. The
\emph{navigator} carries it across these volumes, switching the material at each
boundary, while reusing the single-volume transport unchanged.

\paragraph{World model: one mother, leaf absorbers.}
The volumes are placed in a non-interacting \emph{mother} --- a cylinder of air,
$\rho\le\SI{60}{cm}$ --- that encloses two radially disjoint \emph{daughters}: the
phantom ($\rho\le\SI{8}{cm}$) and the ring ($\SI{38.7}\le\rho\le\SI{42.4}{cm}$). The
daughters are \emph{leaves}: nothing is nested inside either, so a photon inside a leaf
sees only that leaf's own boundary (\S\ref{sec:cyl}--\ref{sec:shell}). The mother is the
only volume that contains others, and it has zero cross section --- a photon never
interacts in air, only traverses it.

\paragraph{Locating a point.}
The material at a point $\mathbf{p}$ is that of the innermost volume containing it ---
daughters first, mother last:
\begin{equation}
  \texttt{locate}(\mathbf{p}) =
  \begin{cases}
    \text{phantom} & \mathbf{p}\in\text{phantom},\\
    \text{ring}    & \mathbf{p}\in\text{ring},\\
    \text{world (air)} & \mathbf{p}\in\text{world},\\
    \varnothing    & \text{otherwise (escaped)}.
  \end{cases}
\end{equation}

\paragraph{Distance to the next material change.}
For the air traverse we need the nearest distance at which membership changes --- leave
the current volume, or enter another:
\begin{equation}
  \texttt{next\_boundary}(\text{current},\mathbf{p},\dvec) =
  \min\Bigl(\,\texttt{exit}(\text{current}),\ \min_{v\neq\text{current}}
  \texttt{entry}(v)\,\Bigr).
\end{equation}
This is well defined because $\texttt{entry}(v)=\infty$ for a $v$ that already contains
$\mathbf{p}$ (\S\ref{sec:cyl}), so the enclosing mother never contributes a spurious
candidate.

\paragraph{The walk.}
The history advances volume by volume. In a \emph{leaf} the photon is transported by the
single-volume routine, which already stops at that leaf's surface and returns the exit
state $(\mathbf{p},\dvec,E)$ and whether the photon left alive; the navigator carries
that state into the next volume. In the \emph{mother} (air, $\Sigma=0$) there is no
interaction to sample, so the step is a straight skip to \texttt{next\_boundary}:
\begin{equation}
  \text{leaf: transport to its surface} \quad\longrightarrow\quad
  \text{air: skip to next\_boundary} \quad\longrightarrow\quad
  \text{leaf: transport \dots}
\end{equation}
until the photon is absorbed in a leaf or the skip carries it out of the world. Treating
air as a skip --- rather than transporting it with the mother's own boundary --- is what
keeps the photon from sailing through the embedded ring to the world edge; it relies on
the mother being non-interacting and the daughters being leaves.

\paragraph{Backscatter across the bore.}
A photon that Compton-scatters in the crystal back toward the patient exits the ring's
\emph{inner} wall --- a leaf exit into the bore (\S\ref{sec:shell}). The bore is air, so
the next skip evaluates \texttt{next\_boundary} from inside the bore, whose
$\texttt{entry}(\text{ring})$ is the \emph{far} wall: the photon re-enters the opposite
crystal and is transported there. The diametral re-entry is produced by the same two
rules, with no special case.

\paragraph{Boundary nudge.}
A photon handed from one volume to the next sits exactly on their shared surface, where
the inclusive ($\le$) membership test is true for both. To re-locate unambiguously the
navigator evaluates \texttt{locate} a nudge $\varepsilon_{\mathrm{nav}}\sim\SI{1}{nm}$
ahead, $\mathbf{p}+\varepsilon_{\mathrm{nav}}\dvec$; the per-solid $t>\texttt{SURFACE\_EPS}$
test (\S\ref{sec:setup}) already discards the surface just crossed.

\end{document}
